% Promoted after complete independent root and companion proof/source/code audits.
\subsection{Prime Heisenberg tensors: dual rank and a birational radical map}

Let \(k\) be algebraically closed of characteristic two, let \(p\) be an
odd prime, and put \(m=(p-1)/2\). Fix a primitive \(p\)-th root
\(\zeta\), and the \(p\)-dimensional Schr\"odinger representation \(W\)
with basis \(e_i\), indexed by \(\mathbf Z/p\).
The shift and clock are \(e_i\mapsto e_{i+1}\) and
\(e_i\mapsto\zeta^i e_i\). Their finite Heisenberg group has centre
\(\boldsymbol\mu_p\), acting on \(W\) with weight one.
Write \(W^{(2)}\) for its Frobenius-twisted representation, whose
central weight is two.

\begin{lemma}[The two actual contraction matrices]
\label{lem:prime-Heisenberg-two-contractions}
\status{Proved; independent audits passed.}
Every equivariant alternating tensor
\(\beta:\bigwedge^2W\to W^{(2)}\) has the form
\[
 \beta(e_i,e_j)=c_{i-j}e_{(i+j)/2}^{(2)},\qquad
 c_0=0,\quad c_a=c_{-a}.                         \tag{PHR1}
\]
Here division by two is in \(\mathbf Z/p\). Its dual contraction
and multiplication by a section are respectively
\[
 B_y(i,j)=c_{i-j}y_{(i+j)/2},\qquad
 M_s(k,j)=c_{2(k-j)}s_{2k-j}.                    \tag{PHR2}
\]
They satisfy \(M_ss=0\), and
\[
                 B_ys=0\quad\Longleftrightarrow\quad
                 M_s^t y=0.                   \tag{PHR3}
\]
\end{lemma}
\begin{proof}
The clock character forces the output index to be \((i+j)/2\);
shift equivariance makes the coefficient depend only on \(i-j\).
Alternation gives exactly the stated conditions on the \(c_a\).
Contracting with \(y\) gives the first matrix. In the second,
the input index multiplying the \(j\)-th basis vector is
\(i=2k-j\), giving the displayed coefficient.
The identity \(M_ss=0\) is \(\beta(s,s)=0\), while (PHR3)
is the equality of the two ways of evaluating
\(y(\beta(s,v))\). In general the second matrix is not an
alternating reindexing of the first; the equality \(2=-1\)
used in the three-dimensional example is special to \(p=3\).
\end{proof}

\begin{proposition}[A complete universal dual-rank certificate through seventeen]
\label{prop:prime-Heisenberg-dual-rank-through-seventeen}
\status{Proved; independent audits passed.}
If
\[
                     p\in\{3,5,7,11,13,17\},
                                                        \tag{PHR4}
\]
then every nonzero coefficient vector \(c\) in (PHR1) has
\(\operatorname{rank}B_y=p-1\) for generic \(y\).
The assertion is over every characteristic-two field; it is
not an inference from a finite set of coefficient values.
\end{proposition}
\begin{proof}
Delete index zero and expand the maximal Pfaffian as a polynomial
over \(\mathbf F_2\):
\[
 P_0(c,y)=
 \sum_{\mathcal M}
 \prod_{\{i,j\}\in\mathcal M}
       c_{\min(i-j,j-i)}y_{(i+j)/2}.             \tag{PHR5}
\]
The sum runs over all perfect matchings of
\(\{1,\ldots,p-1\}\); the representative of each nonzero
distance is chosen in \(\{1,\ldots,m\}\). Repeated monomials
cancel in pairs. A coefficient support is a nonempty subset
\(S\subset\{1,\ldots,m\}\), meaning exactly that \(c_a\ne0\)
for \(a\in S\), and that the other coefficients vanish.

The retained standalone programs reconstruct (PHR5) from all its
matchings. For every support in every prime in (PHR4), the output
specifies a \(y\)-monomial whose coefficient, after setting the
coefficients outside \(S\) to zero, is a single monomial in the
remaining \(c_a\). That coefficient is nonzero for every point
of the indicated support torus. Thus \(P_0(c,y)\) is not the
zero polynomial for any nonzero \(c\).

For clarity, at \(p=5\) the full formula is
\[
 P_0=c_1^2y_1y_4+c_2^2y_2y_3+c_1c_2y_0^2.
                                                        \tag{PHR6}
\]
At \(p=7\), coefficients of actual \(y\)-monomials include
\(c_1c_2^2,c_1^2c_3,c_2c_3^2\). Their vanishing would leave
at most one nonzero \(c_a\); the corresponding pure-cube term
then survives. For the larger primes, the same support argument
is completely enumerated, with an explicit witness for each support.
The support counts are \(1,3,7,31,63,255\), respectively.

For \(p=17\), the second checker reconstructs all
\[
                 15!!=2{,}027{,}025
\]
matchings. After cancellation it obtains \(1{,}465{,}345\)
terms in \(43{,}263\) distinct \(y\)-coefficients, and supplies
all \(255\) support witnesses. It also retains the entire
unrestricted \(c\)-coefficient of each selected witness.
Packed exponents use four bits per variable: every exponent is
at most the total degree eight, so there is no carry between
variables. The smaller-prime program uses literal exponent tuples.
These are exact finite polynomial proofs for the listed primes,
with no claim about an untested prime.
\end{proof}

\begin{lemma}[Prime-degree incompressibility by diagonal trace]
\label{lem:prime-Severi-Brauer-diagonal-trace}
\status{Proved; independent audits passed.}
Let \(D\) be a central division algebra of prime degree \(p\)
over a field \(F\), and let \(X=\operatorname{SB}(D)\).
Every rational self-map \(X\dashrightarrow X\) is dominant.
This statement does not assume separability of the map.
\end{lemma}
\begin{proof}
Put \(n=p-1\). Every closed point of \(X\) has degree divisible
by \(p\). Indeed its residue field \(E\) splits \(D\), so the
left \(D\)-module \(E^p\) has \(F\)-dimension \(p[E:F]\).
As \(D\) is a division algebra of dimension \(p^2\), this
dimension is divisible by \(p^2\). Hence \(p\mid[E:F]\).
The same divisibility holds for the degree of every zero-cycle.

Let \(\Gamma\in\operatorname{CH}^n(X\times X)\) be the closure
of the graph of a rational self-map. Over an algebraic closure,
write
\[
 \overline\Gamma=\sum_{i=0}^{n}a_i h_1^i h_2^{\,n-i},
                 \qquad a_0=1.                 \tag{PHR7}
\]
The last equality follows from the first projection of the graph
being birational. The displayed basis consists of mutually
orthogonal idempotents for composition of correspondences:
the middle-factor intersection has degree one exactly when
the two indices agree. Consequently the split coefficients of
\(\Gamma^{\circ(p-1)}\) are \(a_i^{p-1}\).

Pull this correspondence back by the diagonal regular embedding.
Refined intersection produces a zero-cycle on \(X\), whose
degree after extension of the field is
\[
                         \sum_{i=0}^{p-1}a_i^{p-1}.
                                                        \tag{PHR8}
\]
This integer is divisible by \(p\). Modulo \(p\), each summand
is one if \(p\nmid a_i\), and is zero otherwise. There are
exactly \(p\) summands, and \(a_0=1\); thus all \(p\) of the
coefficients are nonzero modulo \(p\). In particular
\(a_{p-1}\ne0\). This is the multiplicity for the second
projection of the graph, which would be zero if the map were
nondominant. The desired dominance follows.

This is the prime-degree specialization of the diagonal-trace
method underlying Karpenko's incompressibility results; the
complete original proofs of Lemmas~2.21 and~2.23 and Corollaries~2.22 and~2.24
of the 18 April 2009 preprint, and the correspondence argument
in Theorem~2.1 of the October 1999 preprint, were read
\citep{Karpenko2009UpperMotives,Karpenko1999Anisotropy}. The elementary
specialization above supplies the required proof directly in
characteristic two; neither resolution nor a separability
assumption is used.
\end{proof}

\begin{lemma}[A Heisenberg rational map between the two projective spaces is dominant]
\label{lem:prime-Heisenberg-projective-dominance}
\status{Proved; independent audits passed.}
Every Heisenberg-equivariant rational map
\[
            \kappa:\mathbf P((W^{(2)})^*)\dashrightarrow
                              \mathbf P(W)             \tag{PHR9}
\]
is dominant.
\end{lemma}
\begin{proof}
The projective actions factor through
\(D_0=(\mathbf Z/p)^2\). Introduce independent variables \(a,b\),
and work over the infinite field \(F=k(a)((b))\).
Twist by the \(D_0\)-torsor \(F(A,B)/F\), where
\(A^p=a\), \(B^p=b\). Write \(S,C\) for the shift and clock,
so \(CS=\zeta SC\), and let the two Galois generators send
respectively \(A\mapsto\zeta A\) and \(B\mapsto\zeta B\).
The descended matrix algebra is the invariant algebra for the
commuting semilinear actions
\(\operatorname{Ad}(S)\circ\sigma_A\) and
\(\operatorname{Ad}(C)\circ\sigma_B\): it contains
\(X=AC\) and \(Y=B^{-1}S\), with
\[
       X^p=a,\qquad Y^p=b^{-1},\qquad XY=\zeta YX.
\]
Their \(p^2\) monomials are linearly independent over \(F\),
so they form the whole invariant algebra; putting \(u=Y^{-1}\)
gives exactly \(u^p=b\) and \(uX=\zeta Xu\), hence the cyclic
class \(\alpha=(a,b)_p\) used in (PHR10).
For central weight \(r\), the same projective lifts have their
scalar two-cocycle raised to the \(r\)-th power, so the class is
\(r\alpha\). The two twists in (PHR9) therefore have classes
\(-2\alpha\) and \(\alpha\), respectively.

Here \(\alpha\) has index \(p\). A direct verification is useful.
Set \(L=F(a^{1/p})\), with generator
\(\sigma(a^{1/p})=\zeta a^{1/p}\). The cyclic algebra is
\[
 \mathcal D=\bigoplus_{i=0}^{p-1}L u^i,\qquad
       u^p=b,\quad u\ell=\sigma(\ell)u.           \tag{PHR10}
\]
The extension \(L/F\) is unramified of residue degree \(p\)
for the \(b\)-adic valuation. For a nonzero sum in (PHR10),
the numbers \(v_b(\ell_i)+i/p\) have distinct fractional parts,
so there is a unique term of smallest valuation. Under
multiplication the leading coefficient remains nonzero, including
when a power of \(u\) crosses \(p\). Thus \(\mathcal D\) has no
zero divisors, and, being finite dimensional, it is a division
algebra of degree \(p\). This also verifies nontriviality of
the displayed projective cocycle.

The equivariant map descends to a rational map \(X\dashrightarrow
Y\), where \(Y=\operatorname{SB}(\mathcal D)\) and \(X\)
has Brauer class \(-2\alpha\) and dimension \(p-1\).
Over \(F(Y)\), the class \(\alpha\), and hence also \(-2\alpha\),
vanishes. Thus \(X_{F(Y)}\) is a projective space.
Choose an \(F(Y)\)-point in the nonempty open where the descended
\(\kappa\) is defined. Such a point exists because this field is
infinite. It gives a rational map \(Y\dashrightarrow X\) whose
image meets that open, so its composite with \(\kappa\) is a
defined rational self-map of \(Y\).
By the preceding lemma, the composite is dominant. Therefore
\(\kappa\) is dominant after twisting, and consequently also
over the original algebraically closed field \(k\).
This argument explicitly arranges composability; it does not
compose arbitrary rational maps through their indeterminacy loci.
\end{proof}

\begin{theorem}[Dual corank one forces a birational radical map in prime dimension]
\label{thm:prime-Heisenberg-birational-radical}
\status{Proved; independent audits passed.}
For any odd prime \(p\), suppose the nonzero tensor (PHR1)
has generic dual rank \(p-1\). Its maximal-Pfaffian radical map
\[
       \kappa:\mathbf P((W^{(2)})^*)\dashrightarrow
                    \mathbf P(W),\qquad [y]\longmapsto\ker B_y
                                                        \tag{PHR11}
\]
is birational. In particular, \(M_s\) has generic rank \(p-1\).
The maximal Pfaffians have no nonconstant common divisor; the
complement of the dual maximal-rank open has codimension at least two.
All these conclusions therefore hold for every nonzero \(c\)
at every prime in (PHR4).
\end{theorem}
\begin{proof}
On the nonempty dual maximal-rank open \(U\), the one-dimensional
radical defines (PHR11). Equivariance of \(\beta\) makes it
Heisenberg equivariant. The projective dominance lemma applies.
For a fixed \([s]\), the fibre over \([s]\) is
\[
                        U\cap\mathbf P(\ker M_s^t),
                                                        \tag{PHR12}
\]
by (PHR3). Dominance between equal-dimensional varieties makes
this a nonempty zero-dimensional fibre for generic \([s]\).
An open subset of a projective linear space can have that
dimension only if the linear space is a point. Hence
\(\dim\ker M_s^t=1\), or \(\operatorname{rank}M_s=p-1\).
The kernel minors then give a rational inverse to (PHR11).
This proves birationality, including separability, rather than
only generic injectivity on geometric points.

For the last assertion let \(g\) be the greatest common divisor
of the maximal Pfaffians. These have degree \(m<p\). The ideal
they generate is preserved by the Heisenberg action, so uniqueness
of the gcd makes its line \(kg\) a one-dimensional representation.
Every character is trivial on the commutator centre. On the
dual variables that centre has weight \(-2\), so
\(p\mid 2\deg g\). Since \(p\) is odd and
\(0\le\deg g\le m<p\), one has \(\deg g=0\).
Thus there is no divisorial component in their common zero locus.
No height-three assertion is made.
\end{proof}

\begin{remark}[Exact scope, retained computations and sources]
\label{rem:prime-Heisenberg-rank-scope}
\status{Proved; independent audits passed.}
The uniform implication from dual corank one to a birational
radical map holds for every odd prime. The proof of dual corank
one itself is presently restricted to the six primes in (PHR4).
The support witnesses are universal coefficient certificates,
not sampling in a finite field.

The exact programs
\path{code/monodromy/post_CP35_prime_Pfaffian_support_certificates.py}
and
\path{code/monodromy/post_CP35_prime17_Pfaffian_support_certificate.py}
passed at 05:09:26.921919 and 05:13:01.666864 UTC on
20 November 2026, in \(0.03530208900338039\) and
\(2.537690808007028\) seconds. The latter includes the complete
two-million-matching reconstruction and all support witnesses.
A separate direct multiplication-minor reconstruction for
\(p=3,5,7,11\) also passed, at 05:16:54.941064 UTC in
\(1.5080241010000464\) seconds; it is supplementary and is not
used to infer the case seventeen.
All source, outputs and streams are retained.

The source boundary is explicit. Karpenko's original
\emph{Upper motives of algebraic groups and incompressibility of
Severi--Brauer varieties}, 18 April 2009 preprint,
Lemmas~2.21 and~2.23 and Corollaries~2.22 and~2.24, printed pp.~9--10, and
\emph{On anisotropy of orthogonal involutions}, October 1999 preprint,
Theorem~2.1
with its correspondence proof, printed pp.~4--5, were actually
opened and read. The prime-degree argument needed here is
reproved above \citep{Karpenko2009UpperMotives,Karpenko1999Anisotropy}, so no
characteristic restriction from the
orthogonal-involution application is imported.

For a geometric Cartier application one must first identify the
actual theta representation and prove that its tensor is nonzero.
These algebraic results alone do not prove that the Euler
Frobenius is generically bijective, that its wall meets the common
maximal-rank open, that the nonzero Lang cover is connected, or
that the global monodromy group is full. No fixed principal
degree-seventeen group conclusion is asserted.
\end{remark}

% Promotion provenance: frozen candidate and explicit cocycle delta retained.
